\chapter{文明规模的计算动力学：人口、网络与制度能力}
\label{chp:civilization_scaling_law}

讨论文明的规模，不能把人口、道路数量或互联网账户简单替换成模型参数，再把上一章的缩放曲线平移到社会历史。文明是由具有权利、偏好和有限信息的主体，物质基础设施，组织制度，知识存量以及生态约束共同构成的开放系统；同样的人口规模可以对应完全不同的生产、照护、创新与风险分布。我们因此把“文明能力”重建为条件化的多目标表现，而不是单一总量，更不把物理隐喻当作历史必然律。本章将沿着人口与网络的测量、制度化协作、互联网的双重效应、规模陷阱、人机协作和有效多样性六条线索，给出可观察、可干预且允许反例的社会计算动力学。

\section{社会系统的状态基质：人口、关系与制度接口}
\label{sec:social_manifold_substrate}

我们先规定研究对象。设时刻$t$的社会系统由状态
\[
S_t=(P_t,G_t,I_t,K_t,R_t,X_t)
\]
表示。其中$P_t$是主体集合及其人口属性，$G_t=(P_t,E_t,W_t)$是多层关系网络，边可以分别表示家庭照护、交易、运输、通信、知识引用和行政授权，$W_t$记录方向、容量、时延、可靠性与访问资格；$I_t$是法律、组织流程和执行机制构成的制度接口；$K_t$是可被访问、验证和传承的知识与技能存量；$R_t$是土地、能源、设备、财政和劳动时间等资源账本；$X_t$则包含气候、疾病、外部贸易和冲突等环境输入。这里的人口数量$N_t=|P_t|$是计数，边容量可以用人次、吨公里、比特率或货币流量表示，时延用时间表示，资源必须保留各自单位。只有在明确归一化和任务语境后，我们才把这些量组成无量纲指标；人口、信息熵、行政代价和实际焦耳能耗不能写进同一个未说明单位的“社会能量”。

所谓连通度也不是一个充分统计量。平均度相同的两个网络，可能一个由少数枢纽控制，另一个具有多条冗余路径；它们在日常效率上接近，在灾害、审查或供应链中断时却表现相反。对给定任务$q$，我们用可达率$a_q$、端到端时延$\tau_q$、有效吞吐$b_q$、来源可追溯率$v_q$和失效后的恢复时间$r_q$描述网络服务，而不是令$\rho=\langle k\rangle$承担全部解释。道路提高货物可达性，却不必提高政治表达；社交平台扩大消息暴露，也不等于可信知识传播。网络层必须按对象和操作分型，跨层结论需要身份绑定与接口证据。例如，一个研究团队能否把实验结果转化为药物供应，既取决于论文引用网，也取决于试验监管、制造设备、采购合同和基层医疗读出，不能从任一层的“小世界”性质直接推出整体协调。

制度$I_t$在模型中不是抽象的“有效维度”，而是一组改变合法转移集合的规则与实现。我们可把一步社会更新写为
\[
S_{t+1}=\Phi_{\Delta t}(S_t,U_t,\xi_t;I_t,M_t),
\]
其中$U_t$是政策、投资、迁移或集体行动，$\xi_t$是不可控扰动，$M_t$是观测协议，$\Delta t$是离散步长。若法律、组织边界、外部输入或统计口径随时间变化，它们都必须进入更新，不能用固定自治方程掩盖。制度的作用可以表现为降低搜索和协调成本、保护试错、分配责任、提供申诉或保存长期承诺；也可能制造排斥、寻租和单点故障。判断制度是否增加系统能力，需要比较干预前后的任务完成率、分配尾部、错误纠正速度和权利损害，而不是把职业门类或产业数量直接解释为几何维数。

在此基础上，我们把文明表现定义成任务与群体条件下的向量$Y_t=(y_1,\ldots,y_m)$，例如健康寿命、基础物资可靠供给、知识生产、灾害恢复、生态负担和公民可申诉性。聚合函数$F(Y_t;\omega)$中的权重$\omega$属于规范选择，必须公开，不能假装由自然定律唯一给出。人口增加会扩大潜在专门化、市场和样本来源，也会提高住房、照护、协调与生态压力；网络增强会缩短某些路径，也可能放大谣言、金融挤兑和级联故障；制度分工会提高并行处理能力，也可能使局部知识在层级压缩中丢失。因此，本章所说的“规模效应”只是这些相反机制在指定边界和时间窗内的合成结果。

经验检验应至少包含三类对照。第一，在人口相近地区比较不同运输、通信与治理接口，识别连通性而非人口本身的作用；第二，在网络技术近似相同的组织中比较授权、审计与纠错制度，检验制度是否改变可恢复性；第三，追踪自然灾害、线路中断、法规变化或平台迁移等准实验，观察指标是否按模型预测改变。历史案例可以生成机制假说，却不能凭帝国兴衰的事后叙事证明普适阈值。这个状态基质为后文奠定边界：社会结构网络与人口中的分布式知识相互耦合，但二者不是同一对象；任何能力提升都必须通过具体任务、群体后果和可审计接口来验收。
